% Background and Concepts

\subsection{x86-64 Architecture Fundamentals}

The x86-64 architecture (also called AMD64 or Intel 64) extends the 32-bit x86 instruction set with 64-bit addressing and a large set of legacy and modern execution features \cite{intel2023manual}. For this project, the most important behaviors are not normal arithmetic operations but boundary cases where the processor must update architectural state exactly.

\subsubsection{x87 Floating Point Unit}

The x87 FPU provides hardware floating-point arithmetic using an 8-register stack architecture, with registers labeled ST(0) through ST(7). The 16-bit FPU Status Word tracks state such as the Invalid Operation bit, Stack Fault bit, condition-code bits, and the TOP field. When a program pushes a ninth value onto the stack, native hardware sets a specific status-word pattern. Similarly, when a program pops from an empty x87 stack, the processor sets invalid-operation and stack-fault related bits.

These details are easy to simplify in emulators, but they matter because software can observe the exact status word. In the current native baseline, FPU stack overflow returns \texttt{0x3a41}, while empty \texttt{FSTP} returns \texttt{0x0841}. Therefore, x87 tests are strong probes for whether an emulator models more than just ordinary floating-point arithmetic.

\subsubsection{Flags, Faults, and Timestamp Instructions}

The EFLAGS register contains CPU status and control flags. The \texttt{LAHF} instruction copies low flag bits into AH, including the Auxiliary Flag and reserved bits. Because the expected native value in the current probe is compact (\texttt{0x03}), even a small discrepancy becomes an effective fingerprint.

Exception behavior is equally important. Signed division overflow through \texttt{IDIV} should raise \texttt{SIGFPE}, while an aligned SSE load such as \texttt{MOVDQA} should fault when pointed at unaligned memory. These are high-confidence correctness probes because the expected behavior is architectural and not merely environment-specific.

Timestamp behavior is useful in a different way. \texttt{RDTSC} should be monotonic across immediate consecutive reads, while \texttt{RDTSCP} also returns a processor identifier through TSC\_AUX. Exact TSC\_AUX values can vary across processors and kernels, but consistent synthetic values still provide useful emulator-detection signals.

\subsubsection{SSE Rounding and Machine-State Instructions}

The MXCSR register controls SSE floating-point behavior, including rounding mode. A probe that requests round-down and then converts a floating-point value checks whether the emulator honors this control state. In the current native run, the expected result is \texttt{1}; different results indicate a semantic gap in SSE state handling.

The system-state instructions \texttt{SGDT}, \texttt{SIDT}, \texttt{SLDT}, \texttt{STR}, and \texttt{SMSW} expose descriptor-table or machine-status information. Their exact values are environment-sensitive, but they are valuable fingerprints because emulators often fault, return zero, or synthesize values instead of matching host-native state. Consequently, this report treats arithmetic, fault, and status-word mismatches as stronger correctness evidence, while treating machine-state differences primarily as emulator-identification evidence.

\subsection{Emulation Techniques}

Modern CPU emulators use different execution models. QEMU TCG translates guest instructions to an intermediate representation before generating host-native code \cite{bellard2005qemu}. Blink focuses on compact x86-64 binary execution, Box64 targets x86-64 user programs on non-x86 hosts, and Unicorn exposes a library interface derived from QEMU for controlled instruction execution \cite{unicorn2015}. Other tools such as Icicle, MWEMU, and KUBERA provide shellcode-oriented or analysis-oriented execution paths.

Each model creates different comparison constraints. Full Linux user-mode execution can be compared through process behavior and signal delivery, while shellcode-only execution often uses emulator-specific wrappers and return conventions. Consequently, this project separates process-backed probes from shellcode-backed probes instead of pretending that all backends expose identical evidence.
